% =====================================================================
%  Kiến trúc máy tính RISC-V 32-bit (RV32I) trên Logisim Evolution
%  Tài liệu học từng bước, từng khối, có hình TikZ.
%  Biên dịch:  pdflatex KienTruc_RV32I.tex   (chạy 2 lần để có mục lục)
%  Hoặc dùng Overleaf (compiler: pdfLaTeX).
% =====================================================================
\documentclass[11pt,a4paper]{article}

\usepackage[utf8]{inputenc}
\usepackage[T5]{fontenc}
\usepackage{lmodern}
\usepackage[vietnamese]{babel}
\usepackage[a4paper,margin=2.2cm]{geometry}
\usepackage{amsmath,amssymb}
\usepackage{booktabs,longtable,array,enumitem,tabularx}
\usepackage{microtype}
\usepackage{xcolor}
\usepackage{tikz}
\usepackage[most]{tcolorbox}
\usetikzlibrary{arrows.meta,positioning,shapes.geometric,calc,fit,backgrounds}
\usepackage[unicode,colorlinks=true,linkcolor=blue!50!black,urlcolor=blue!50!black]{hyperref}

\setlist{itemsep=2pt,topsep=3pt,parsep=0pt}
\setlength{\parindent}{0pt}
\setlength{\parskip}{4pt}
\renewcommand{\arraystretch}{1.15}
\setlength{\emergencystretch}{3em}

% tên tín hiệu / lệnh (chữ đơn cách, gạch dưới an toàn)
\newcommand{\sig}[1]{\texttt{\detokenize{#1}}}

% hộp "Thí nghiệm", "Câu hỏi tự kiểm tra", "Lưu ý"
\newtcolorbox{khung}[2]{enhanced,breakable,colback=#1!6,colframe=#1!70!black,
  boxrule=0.6pt,arc=2pt,left=5pt,right=5pt,top=3pt,bottom=3pt,
  fonttitle=\bfseries\small,coltitle=black,colbacktitle=#1!18,title={#2},
  before skip=6pt,after skip=6pt}
\newenvironment{thinghiem}{\begin{khung}{green}{Thí nghiệm}}{\end{khung}}
\newenvironment{kiemtra}{\begin{khung}{orange}{Câu hỏi tự kiểm tra}}{\end{khung}}
\newenvironment{luuy}{\begin{khung}{red}{Lưu ý}}{\end{khung}}

% kiểu TikZ dùng chung
\tikzset{
  blk/.style={draw,rounded corners=2pt,minimum height=8mm,minimum width=16mm,
              align=center,font=\footnotesize,fill=blue!6},
  ctl/.style={blk,fill=orange!15},
  mem/.style={blk,fill=green!10},
  reg/.style={blk,fill=yellow!20,minimum width=4mm,minimum height=26mm},
  % MUX: hình thang xoay 270 độ, nên "minimum width" là chiều CAO khi vẽ ra
  mux/.style={draw,trapezium,trapezium angle=75,trapezium stretches=true,shape border rotate=270,
              minimum width=12mm,minimum height=7mm,inner sep=1pt,align=center,font=\scriptsize,fill=gray!12},
  ar/.style={-{Stealth[length=2mm]},thick},
  bus/.style={-{Stealth[length=2mm]},line width=1.2pt},
  cs/.style={-{Stealth[length=2mm]},thick,orange!70!black,dashed},
  lab/.style={font=\scriptsize,inner sep=1pt},
  stg/.style={draw,minimum width=14mm,minimum height=5mm,font=\scriptsize,inner sep=1pt}
}

\title{\textbf{Kiến trúc máy tính RISC-V 32-bit (RV32I)\\ trên Logisim Evolution}\\[4pt]
\large Tài liệu học từng bước: từng khối, bên trong có gì, vì sao}
\author{Nhóm bài tập lớn --- bản \LaTeX{} tiếng Việt}
\date{}

\begin{document}
\maketitle
\tableofcontents
\newpage

% =====================================================================
\section{Mục tiêu và cách dùng tài liệu}
% =====================================================================
Sau khi học hết tài liệu này, bạn có thể giải thích \textbf{mọi circuit của dự án, bên trong có gì và vì sao}, từ một cổng logic đến CPU pipeline 5 tầng, và chứng minh bằng kiểm thử rằng nó chạy đúng.

\textbf{Dự án gồm:} hai CPU dùng chung các khối: \sig{main} (single-cycle) và \sig{Pipeline} (5 tầng IF, ID, EX, MEM, WB); 14 circuit; bộ nhớ kiểu Harvard (ROM lệnh 256 từ, RAM dữ liệu 256 từ); assembler, bộ mô phỏng tham chiếu và bộ kiểm thử tự động.

\textbf{Cách học:} làm lần lượt các bước. Mỗi bước gồm \emph{đọc}, \emph{nhìn} (mở circuit trong Logisim), \emph{làm} (một thí nghiệm nhỏ) và \emph{kiểm tra} (trả lời không nhìn tài liệu). Chưa trả lời được thì lặp lại bước đó.

\begin{table}[htbp]
\centering
\small
\begin{tabularx}{\linewidth}{@{}cc>{\raggedright}p{4.6cm}>{\raggedright\arraybackslash}X@{}}
\toprule
\textbf{Giai đoạn} & \textbf{Buổi} & \textbf{Chủ đề} & \textbf{Circuit} \\
\midrule
0 & 1 & Công cụ và bức tranh lớn & (nhìn từ ngoài) \\
1 & 2--3 & Nền tảng: CPU chạy chương trình, định dạng lệnh & --- \\
2 & 4 & Giải mã và điều khiển & \sig{InstructionDecode}, \sig{MainControl}, \sig{ImmGen} \\
3 & 5 & Thanh ghi và số học & \sig{RegisterFile}, \sig{ALUControl}, \sig{ALU}, \sig{BranchUnit} \\
4 & 6--7 & Bộ nhớ và CPU single-cycle & \sig{Memory}, \sig{LoadExtract}, \sig{StoreMerge}, \sig{main} \\
5 & 8--10 & Pipeline & \sig{Pipeline}, \sig{ForwardingUnit}, \sig{HazardUnit} \\
6 & 11--12 & Chương trình, kiểm thử, công cụ, demo & assembler, simulator, test \\
\bottomrule
\end{tabularx}
\caption{Lộ trình 12 buổi (mỗi buổi khoảng 90 phút, tổng cộng khoảng 18 giờ).}
\end{table}

Nếu ba người chia việc: mỗi người học sâu một phần (xem bảng ở phần~\ref{sec:chia}) nhưng cả ba cùng làm giai đoạn 0, 1 và 6.

% =====================================================================
\section{Giai đoạn 0 --- Công cụ và bức tranh lớn (buổi 1)}
% =====================================================================
\subsection*{Bước 0.1: mở dự án}
\begin{enumerate}
  \item Cần Java 21. Mở dự án bằng lệnh:\\ \texttt{java -jar logisim-evolution-4.1.0-all.jar RiscVComputer.circ}
  \item Trong cây bên trái có 14 circuit. Chia thành ba nhóm:
  \begin{itemize}
    \item hai CPU: \sig{main}, \sig{Pipeline};
    \item các khối datapath: \sig{InstructionDecode}, \sig{MainControl}, \sig{ImmGen}, \sig{RegisterFile}, \sig{ALUControl}, \sig{ALU}, \sig{BranchUnit}, \sig{Memory}, \sig{LoadExtract}, \sig{StoreMerge};
    \item các khối chỉ có trong pipeline: \sig{ForwardingUnit}, \sig{HazardUnit}.
  \end{itemize}
\end{enumerate}

\subsection*{Bước 0.2: các thao tác Logisim cần biết}
\begin{center}
\begin{tabular}{@{}ll@{}}
\toprule
\textbf{Phím / thao tác} & \textbf{Tác dụng} \\
\midrule
Ctrl+R & đặt lại mô phỏng (PC và mọi thanh ghi về 0) \\
Ctrl+F9 & một chu kỳ xung nhịp đầy đủ (một tick) \\
Tab State & hiện \sig{PC} và \sig{x0..x31} trực tiếp \\
Tunnel (đường hầm) & các tunnel cùng tên là \textbf{cùng một dây} \\
Hàng probe cuối mạch & các giá trị cần theo dõi khi demo \\
\bottomrule
\end{tabular}
\end{center}

\subsection*{Bước 0.3: bức tranh lớn}
Hình~\ref{fig:big} cho thấy ý tưởng chung: CPU lấy lệnh từ ROM, giải mã, thực thi bằng ALU, truy cập RAM và ghi kết quả vào thanh ghi.

\begin{figure}[htbp]
\centering
\begin{tikzpicture}[node distance=9mm]
  \node[mem] (rom) {ROM lệnh\\256$\times$32};
  \node[blk,right=of rom] (dec) {Giải mã\\+ điều khiển};
  \node[blk,right=of dec] (rf) {Thanh ghi\\x0..x31};
  \node[blk,right=of rf] (alu) {ALU};
  \node[mem,right=of alu] (ram) {RAM dữ liệu\\256$\times$32};
  \node[blk,above=10mm of rom] (pc) {PC};
  \draw[ar] (pc) -- (rom);
  \draw[ar] (rom) -- node[lab,above]{lệnh} (dec);
  \draw[ar] (dec) -- (rf);
  \draw[ar] (rf) -- (alu);
  \draw[ar] (alu) -- (ram);
  \draw[ar] (alu.north) |- ++(0,6mm) -| (pc.east) node[lab,pos=.25,above]{địa chỉ nhảy / PC+4};
  \draw[ar] (ram.south) -- ++(0,-7mm) -| node[lab,pos=.25,below]{ghi kết quả (write-back)} (rf.south);
\end{tikzpicture}
\caption{Bức tranh lớn: nạp lệnh $\rightarrow$ giải mã $\rightarrow$ thanh ghi $\rightarrow$ ALU $\rightarrow$ bộ nhớ $\rightarrow$ ghi lại.}
\label{fig:big}
\end{figure}

\begin{kiemtra}
\begin{itemize}
  \item \sig{main} khác \sig{Pipeline} ở điểm nào?
  \item Vì sao có hai bộ nhớ (ROM lệnh, RAM dữ liệu)?
  \item Tunnel dùng để làm gì?
\end{itemize}
\end{kiemtra}

% =====================================================================
\section{Giai đoạn 1 --- Nền tảng (buổi 2--3)}
% =====================================================================
\subsection{CPU chạy chương trình như thế nào}
\begin{enumerate}
  \item \textbf{Nạp lệnh (fetch):} đọc lệnh ở địa chỉ trong \sig{PC}.
  \item \textbf{Giải mã (decode):} tách trường lệnh, đọc thanh ghi, sinh tín hiệu điều khiển.
  \item \textbf{Thực thi (execute):} ALU tính; quyết định rẽ nhánh.
  \item \textbf{Bộ nhớ (memory):} load / store.
  \item \textbf{Ghi lại (write-back):} đưa kết quả vào thanh ghi \sig{rd}; \sig{PC} chuyển sang lệnh kế tiếp.
\end{enumerate}
Các ý quan trọng: trạng thái chỉ đổi ở \textbf{cạnh xung nhịp}; mạch tổ hợp khác mạch tuần tự; số âm dùng bù hai (ví dụ $-1 = \texttt{0xFFFFFFFF}$); mọi giá trị trong tài liệu viết ở hệ 16 (hex) nếu có tiền tố \texttt{0x}.

\begin{thinghiem}
Chạy tay \texttt{programs/demo.s} 3 vòng lặp trên giấy (PC, x10, x11) rồi so với bảng 40 tick ở phần~\ref{sec:demo}.
\end{thinghiem}

\subsection{Định dạng lệnh RISC-V}
Có sáu định dạng R, I, S, B, U, J. Các trường nằm \textbf{cố định} ở cùng vị trí bit nên giải mã chỉ là đi dây.

\begin{center}
\small
\setlength{\tabcolsep}{3pt}
\begin{tabular}{|c|c|c|c|c|c|c|}
\hline
\textbf{Kiểu} & \textbf{31--25} & \textbf{24--20} & \textbf{19--15} & \textbf{14--12} & \textbf{11--7} & \textbf{6--0}\\
\hline
R & funct7 & rs2 & rs1 & funct3 & rd & opcode\\
\hline
I & \multicolumn{2}{c|}{imm[11:0]} & rs1 & funct3 & rd & opcode\\
\hline
S & imm[11:5] & rs2 & rs1 & funct3 & imm[4:0] & opcode\\
\hline
B & imm[12$|$10:5] & rs2 & rs1 & funct3 & imm[4:1$|$11] & opcode\\
\hline
U & \multicolumn{4}{c|}{imm[31:12]} & rd & opcode\\
\hline
J & \multicolumn{4}{c|}{imm[20$|$10:1$|$11$|$19:12]} & rd & opcode\\
\hline
\end{tabular}
\end{center}

\textbf{Chín opcode đang dùng:}
\texttt{0x33} (R), \texttt{0x13} (I-ALU), \texttt{0x03} (load), \texttt{0x23} (store), \texttt{0x63} (branch), \texttt{0x37} (LUI), \texttt{0x17} (AUIPC), \texttt{0x6F} (JAL), \texttt{0x67} (JALR).

\begin{thinghiem}
Chạy \texttt{python tools/rv.py asm programs/demo.s -l}, rồi giải mã tay ba từ lệnh \texttt{00b50533}, \texttt{fff58593}, \texttt{fe059ce3}: viết các trường dạng nhị phân, nêu định dạng và tên lệnh.
\end{thinghiem}
\begin{kiemtra}
\begin{itemize}
  \item Định dạng nào không có \sig{rd}?
  \item Vì sao hằng số kiểu B và J thiếu bit 0?
  \item \sig{funct7[5]} làm thay đổi điều gì?
\end{itemize}
\end{kiemtra}

% =====================================================================
\section{Giai đoạn 2 --- Giải mã và điều khiển (buổi 4)}
% =====================================================================

\subsection{Khối \texorpdfstring{\texttt{InstructionDecode}}{InstructionDecode}}
\begin{itemize}
  \item \textbf{Mục đích:} cắt lệnh 32 bit thành các trường.
  \item \textbf{Vào / ra:} vào \sig{Instruction}; ra \sig{opcode}(7), \sig{rd}(5), \sig{funct3}(3), \sig{rs1}(5), \sig{rs2}(5), \sig{funct7}(7).
  \item \textbf{Bên trong:} đúng một \emph{splitter} (32 $\rightarrow$ 6 phần). Không có cổng logic.
  \item \textbf{Vì sao:} RISC-V đặt số hiệu thanh ghi luôn ở cùng chỗ, nên có thể đọc thanh ghi \emph{trong khi} còn đang giải mã opcode.
\end{itemize}

\begin{figure}[htbp]
\centering
\begin{tikzpicture}[x=1mm,y=1mm]
  \draw (0,0) rectangle (112,8);
  \foreach \x/\t in {0/{funct7 [31:25]},28/{rs2 [24:20]},48/{rs1 [19:15]},68/{funct3 [14:12]},88/{rd [11:7]}} {
    \draw (\x,0) -- (\x,8);
    \node[font=\tiny,anchor=west] at (\x+1,4) {\t};}
  \node[font=\tiny,anchor=west] at (106,4) {};
  \draw (104,0) -- (104,8);
  \node[font=\tiny] at (108,4) {op};
  \foreach \x in {14,38,58,78,96,108} \draw[ar] (\x,0) -- (\x,-7);
  \node[lab,below] at (14,-7) {\sig{funct7}};
  \node[lab,below] at (38,-7) {\sig{rs2}};
  \node[lab,below] at (58,-7) {\sig{rs1}};
  \node[lab,below] at (78,-7) {\sig{funct3}};
  \node[lab,below] at (96,-7) {\sig{rd}};
  \node[lab,below] at (108,-7) {\sig{opcode}};
\end{tikzpicture}
\caption{\sig{InstructionDecode}: chỉ là đi dây từ các trường bit cố định.}
\end{figure}

\begin{thinghiem}
Gán \sig{Instruction} = \texttt{00b50533} rồi đọc sáu ngõ ra ($\mathtt{opcode}=\mathtt{0x33}$, $\mathtt{rd}=10$, $\mathtt{rs1}=10$, $\mathtt{rs2}=11$).
\end{thinghiem}
\begin{kiemtra}Các bit nào tạo \sig{rs2}? Vì sao khối này không có cổng logic nào?\end{kiemtra}

\subsection{Khối \texorpdfstring{\texttt{MainControl}}{MainControl}}
\begin{itemize}
  \item \textbf{Mục đích:} từ \sig{opcode} sinh tất cả tín hiệu điều khiển của datapath.
  \item \textbf{Vào / ra:} vào \sig{opcode}(7); ra \sig{RegWrite}, \sig{ALUSrc}, \sig{MemWrite}, \sig{MemRead}, \sig{IsBranch}, \sig{IsJAL}, \sig{IsJALR}, \sig{ALUOp}(2), \sig{ImmSel}(3), \sig{ASel}(2), \sig{WBSel}(2).
  \item \textbf{Bên trong:} chín \emph{bộ so sánh} 7 bit (mỗi opcode một cái, kèm \emph{hằng số} 7 bit); các \emph{cổng OR} (tối đa 3 ngõ vào) cho biết lớp lệnh nào bật tín hiệu nào; các \emph{splitter} ghép bit đơn thành bus 2--3 bit.
\end{itemize}

\textbf{Các phương trình ghép bus} (viết $\{a,b,\dots\}$: phần tử đầu là bit cao nhất):
\begin{align*}
\mathtt{ALUOp} &= \{\mathtt{IsR}\lor\mathtt{IsI},\ \mathtt{IsI}\lor\mathtt{IsBranch}\}\\
\mathtt{ImmSel} &= \{\mathtt{IsJAL},\ \mathtt{IsBranch}\lor\mathtt{IsLui}\lor\mathtt{IsAuipc},\ \mathtt{IsStore}\lor\mathtt{IsLui}\lor\mathtt{IsAuipc}\}\\
\mathtt{ASel} &= \{\mathtt{IsLui},\ \mathtt{IsAuipc}\}\\
\mathtt{WBSel} &= \{\mathtt{IsJAL}\lor\mathtt{IsJALR},\ \mathtt{IsLoad}\}
\end{align*}
Ý nghĩa các mã:
\begin{itemize}
  \item \sig{ALUOp}: 00 cộng, 01 trừ, 10 tra bảng kiểu R, 11 tra bảng kiểu I;
  \item \sig{ImmSel}: 000 I, 001 S, 010 B, 011 U, 100 J;
  \item \sig{ASel} (toán hạng A): 00 \sig{rs1}, 01 \sig{PC} (AUIPC), 10 hằng 0 (LUI);
  \item \sig{WBSel} (dữ liệu ghi lại): 00 kết quả ALU, 01 dữ liệu bộ nhớ, 10 \sig{PC+4} (JAL, JALR).
\end{itemize}
Opcode lạ làm mọi tín hiệu bằng 0 (một lệnh NOP vô hại).

\begin{figure}[htbp]
\centering
\begin{tikzpicture}[node distance=7mm]
  \node[lab] (op) {\sig{opcode}};
  \node[blk,right=10mm of op,minimum height=24mm] (cmp) {9 bộ\\so sánh\\7 bit};
  \node[blk,right=26mm of cmp,minimum height=24mm] (or) {cổng OR\\+ splitter};
  \node[lab,right=12mm of or,align=left] (out) {\sig{RegWrite} \sig{ALUSrc}\\\sig{MemWrite} \sig{MemRead}\\\sig{IsBranch} \sig{IsJAL} \sig{IsJALR}\\\sig{ALUOp} \sig{ImmSel}\\\sig{ASel} \sig{WBSel}};
  \draw[bus] (op) -- (cmp);
  \draw[ar] (cmp) -- node[lab,above]{9 cờ lớp lệnh} (or);
  \draw[bus] (or) -- (out);
\end{tikzpicture}
\caption{\sig{MainControl}: so sánh opcode $\rightarrow$ cờ lớp lệnh $\rightarrow$ OR $\rightarrow$ tín hiệu điều khiển.}
\end{figure}

\begin{thinghiem}
Điền bảng tín hiệu điều khiển trên giấy cho \texttt{lw}, \texttt{sw}, \texttt{beq}, \texttt{lui}, \texttt{jal}, rồi đối chiếu với \texttt{tests/maincontrol.vec}.
\end{thinghiem}
\begin{kiemtra}Với \texttt{sw}, những tín hiệu nào bằng 1? \sig{ASel}=10 nghĩa là gì và dùng cho lệnh nào? Chuyện gì xảy ra với opcode 0?\end{kiemtra}

\subsection{Khối \texorpdfstring{\texttt{ImmGen}}{ImmGen}}
\begin{itemize}
  \item \textbf{Mục đích:} dựng lại hằng số 32 bit theo định dạng do \sig{ImmSel} chọn.
  \item \textbf{Vào / ra:} vào \sig{Instruction}, \sig{ImmSel}(3); ra \sig{Imm}(32).
  \item \textbf{Bên trong:} 9 \emph{splitter} (cắt các trường rải rác và ghép đúng thứ tự, trong đó một splitter 32 bit ghép 12 bit 0 với \sig{instr[31:12]} cho kiểu U); 4 \emph{bộ mở rộng bit} mở rộng dấu cho I, S, B, J (12, 12, 13, 21 bit $\rightarrow$ 32); với kiểu U là hằng số 12 bit 0 và splitter ghép (\texttt{imm}$\ll 12$); hai hằng số 1 bit 0 cho bit 0 luôn bằng 0 của B và J; một \emph{multiplexer} 8-ra-1 (dùng 5 ngõ) chọn theo \sig{ImmSel} (0 I, 1 S, 2 B, 3 U, 4 J).
  \item \textbf{Vì sao:} các bit hằng số bị xáo trộn có chủ đích để đa số bit nằm cùng vị trí giữa các định dạng; splitter xếp lại; mở rộng dấu giữ số âm vẫn âm.
\end{itemize}

\begin{figure}[htbp]
\centering
\begin{tikzpicture}[node distance=7mm]
  \node[lab] (in) {\sig{Instruction}};
  \node[blk,right=8mm of in,yshift=12mm] (i) {I: 12 bit\\mở rộng dấu};
  \node[blk,below=2mm of i] (s) {S: 12 bit\\mở rộng dấu};
  \node[blk,below=2mm of s] (b) {B: 13 bit\\mở rộng dấu};
  \node[blk,below=2mm of b] (u) {U: $\mathtt{imm}\ll12$};
  \node[blk,below=2mm of u] (j) {J: 21 bit\\mở rộng dấu};
  \node[mux,right=18mm of s,minimum width=40mm,minimum height=9mm,yshift=-8mm] (m) {MUX\\8:1};
  \node[lab,right=8mm of m] (out) {\sig{Imm}};
  \node[ctl,below=6mm of m] (sel) {\sig{ImmSel}};
  \foreach \n in {i,s,b,u,j} {\draw[ar] (in.east) -- (\n.west); \draw[ar] (\n.east) -- (m.west);}
  \draw[bus] (m) -- (out);
  \draw[cs] (sel) -- (m);
\end{tikzpicture}
\caption{\sig{ImmGen}: năm cách dựng hằng số, multiplexer chọn theo \sig{ImmSel}.}
\end{figure}

\begin{thinghiem}
Giải mã \texttt{fe059ce3} (kiểu B): hằng số phải là \texttt{0xfffffff8} $=-8$. Kiểm tra trong \texttt{tests/immgen.vec}.
\end{thinghiem}
\begin{kiemtra}Vì sao độ lệch của B luôn chẵn? Bit dấu lấy từ đâu ở mỗi định dạng?\end{kiemtra}

\textbf{Bài kiểm tra cuối giai đoạn:} từ một từ lệnh, nêu định dạng, ghi \sig{rd}, \sig{rs1}, \sig{rs2}, hằng số và các tín hiệu điều khiển mà không nhìn tài liệu.

% =====================================================================
\section{Giai đoạn 3 --- Thanh ghi và số học (buổi 5)}
% =====================================================================

\subsection{Khối \texorpdfstring{\texttt{RegisterFile}}{RegisterFile}}
\begin{itemize}
  \item \textbf{Mục đích:} 32 thanh ghi $\times$ 32 bit, mỗi chu kỳ đọc hai và ghi một.
  \item \textbf{Vào / ra:} vào \sig{write_data}, \sig{CLK}, \sig{write} (\sig{RegWrite}), \sig{rd}, \sig{rs1}, \sig{rs2}; ra \sig{rs1_data}, \sig{rs2_data}.
  \item \textbf{Bên trong:} 32 \emph{thanh ghi} (nhãn \sig{x0..x31}, sườn lên); một \emph{demultiplexer} (chọn 5 bit, bit dữ liệu = \sig{RegWrite}) dẫn tín hiệu cho phép ghi tới đúng thanh ghi \sig{rd}; hai \emph{multiplexer} 32-ra-1 làm cổng đọc (bất đồng bộ); hằng số 0 nối vào \sig{x0}.
  \item \textbf{Cách hoạt động:} ghi: chỉ thanh ghi số \sig{rd} có WE = 1; đọc: multiplexer luôn hiện thanh ghi do \sig{rs1}/\sig{rs2} chọn. \sig{x0} chỉ nạp được 0 nên luôn bằng 0.
  \item \textbf{Vì sao:} đọc bất đồng bộ để CPU single-cycle có toán hạng ngay trong chu kỳ; trong pipeline khối này dùng \textbf{xung nhịp đảo} (ghi ở sườn xuống).
\end{itemize}

\begin{figure}[htbp]
\centering
\begin{tikzpicture}[node distance=6mm]
  \node[blk] (demux) {DEMUX\\1$\rightarrow$32\\(chọn \sig{rd})};
  \node[blk,right=26mm of demux,minimum height=30mm] (regs) {32 thanh ghi\\\sig{x0}..\sig{x31}\\(x0 luôn 0)};
  \node[mux,right=14mm of regs,yshift=8mm,minimum width=12mm,minimum height=8mm] (m1) {MUX\\32:1};
  \node[mux,right=14mm of regs,yshift=-8mm,minimum width=12mm,minimum height=8mm] (m2) {MUX\\32:1};
  \node[lab,right=8mm of m1] {\sig{rs1_data}};
  \node[lab,right=8mm of m2] {\sig{rs2_data}};
  \node[lab,left=8mm of demux,yshift=3mm] (rw) {\sig{RegWrite}};
  \node[lab,left=8mm of demux,yshift=-3mm] (rd) {\sig{rd}};
  \draw[ar] (rw) -- (demux); \draw[ar] (rd) -- (demux);
  \draw[ar] (demux) -- node[lab,above,align=center]{32 tín hiệu\\WE} (regs);
  \draw[ar] (regs.east) -- (m1.west); \draw[ar] (regs.east) -- (m2.west);
  \node[ctl,above=4mm of m1] (r1) {\sig{rs1}}; \draw[cs] (r1) -- (m1);
  \node[ctl,below=4mm of m2] (r2) {\sig{rs2}}; \draw[cs] (r2) -- (m2);
  \node[lab,above=6mm of regs] (wd) {\sig{write_data}}; \draw[ar] (wd) -- (regs);
\end{tikzpicture}
\caption{\sig{RegisterFile}: demux dẫn WE tới \sig{rd}; hai multiplexer là hai cổng đọc.}
\end{figure}

\begin{thinghiem}Trong \sig{main}, bấm 3 tick của demo rồi mở \sig{RegisterFile}: tìm thanh ghi nào vừa đổi và đầu ra demux nào bằng 1.\end{thinghiem}
\begin{kiemtra}Điều gì xảy ra khi $\mathtt{rd}=0$ và $\mathtt{RegWrite}=1$? Vì sao đọc bất đồng bộ còn ghi theo xung nhịp?\end{kiemtra}

\subsection{Khối \texorpdfstring{\texttt{ALUControl}}{ALUControl}}
\begin{itemize}
  \item \textbf{Mục đích:} biến \sig{ALUOp}, \sig{funct3}, \sig{funct7} thành mã 4 bit \sig{ALUCtrl}.
  \item \textbf{Vào / ra:} vào \sig{funct7}(7), \sig{ALUOp}(2), \sig{funct3}(3); ra \sig{ALUCtrl}(4).
  \item \textbf{Bên trong:} 19 hằng số 4 bit (các mã: 0 ADD, 1 SUB, 2 AND, 3 OR, 4 XOR, 5 SLL, 6 SRL, 7 SRA, 8 SLT, 9 SLTU) và 5 multiplexer: hai 8-ra-1 (bảng kiểu R và kiểu I, chọn theo \sig{funct3}), hai 2-ra-1 (ADD/SUB và SRL/SRA, chọn theo \sig{funct7[5]}), một 4-ra-1 (chọn theo \sig{ALUOp}: ép ADD, ép SUB, bảng R, bảng I).
  \item \textbf{Vì sao:} một bảng mã = mỗi dòng một hằng số, multiplexer chọn dòng; điều khiển hai tầng (\sig{MainControl} $\rightarrow$ \sig{ALUOp} $\rightarrow$ \sig{ALUControl}) giữ các bảng nhỏ.
\end{itemize}

\begin{table}[htbp]
\centering
\small
\begin{tabularx}{\linewidth}{@{}c>{\raggedright}p{3.4cm}>{\raggedright}p{3.2cm}>{\raggedright\arraybackslash}X@{}}
\toprule
\textbf{\sig{ALUOp}} & \textbf{Nhóm lệnh} & \textbf{\sig{ALUCtrl}} & \textbf{Ghi chú}\\
\midrule
00 & load, store, LUI, AUIPC, JALR & ADD & luôn cộng\\
01 & branch & SUB & so sánh bằng phép trừ\\
10 & kiểu R & theo \sig{funct3}, \sig{funct7[5]} & ADD/SUB, SRL/SRA phân biệt bằng \sig{funct7[5]}\\
11 & kiểu I & theo \sig{funct3} & chỉ lệnh dịch xét \sig{funct7[5]}\\
\bottomrule
\end{tabularx}
\caption{Hai tầng điều khiển ALU.}
\end{table}

\begin{thinghiem}Viết trên giấy mã của \texttt{sub}, \texttt{sra}, \texttt{srai}, \texttt{slti}, \texttt{lw}, \texttt{beq} rồi đối chiếu \texttt{tests/alucontrol.vec}.\end{thinghiem}
\begin{kiemtra}Vì sao \texttt{addi} bỏ qua \sig{funct7} còn \texttt{srai} thì không? Load và branch dùng \sig{ALUOp} nào, vì sao?\end{kiemtra}

\subsection{Khối \texorpdfstring{\texttt{ALU}}{ALU}}
\begin{itemize}
  \item \textbf{Mục đích:} tính kết quả theo mã \sig{ALUCtrl}.
  \item \textbf{Vào / ra:} vào \sig{rs1_data}, \sig{rs2_data}, \sig{ALU_Ctrl}; ra \sig{result}.
  \item \textbf{Bên trong:} bộ cộng, bộ trừ, cổng AND/OR/XOR 32 bit, ba bộ dịch (logic trái, logic phải, số học phải; lượng dịch \sig{B[4:0]}), bộ so sánh có dấu và không dấu, mỗi bộ theo sau là bộ mở rộng 0 từ 1 lên 32 bit (SLT/SLTU trả 0 hoặc 1), một multiplexer 16-ra-1 (dùng 10 ngõ).
  \item \textbf{Cách hoạt động:} mọi phép tính chạy song song, multiplexer chọn một kết quả theo \sig{ALUCtrl}.
\end{itemize}

\begin{figure}[htbp]
\centering
\begin{tikzpicture}[node distance=3mm]
  \node[lab] (a) {\sig{rs1_data}\ (A)};
  \node[lab,below=14mm of a] (b) {\sig{rs2_data}\ (B)};
  \node[blk,right=18mm of a,yshift=18mm] (add) {ADD};
  \node[blk,below=2mm of add] (sub) {SUB};
  \node[blk,below=2mm of sub] (log) {AND / OR / XOR};
  \node[blk,below=2mm of log] (sh) {SLL / SRL / SRA};
  \node[blk,below=2mm of sh] (cmp) {SLT / SLTU\\(so sánh + mở rộng 0)};
  \node[mux,right=22mm of log,minimum width=50mm,minimum height=10mm] (m) {MUX\\16:1};
  \node[lab,right=8mm of m] (res) {\sig{result}};
  \node[ctl,above=5mm of m] (c) {\sig{ALUCtrl}};
  \foreach \n in {add,sub,log,sh,cmp} {\draw[ar] ($(a)!0.5!(b)+(10mm,0)$) -- (\n.west); \draw[ar] (\n.east) -- (m.west);}
  \draw[ar] (a) -- ($(a)!0.5!(b)+(10mm,0)$);
  \draw[ar] (b) -- ($(a)!0.5!(b)+(10mm,0)$);
  \draw[bus] (m) -- (res);
  \draw[cs] (c) -- (m);
\end{tikzpicture}
\caption{\sig{ALU}: tính song song mọi phép, multiplexer chọn theo \sig{ALUCtrl}.}
\end{figure}

\begin{thinghiem}Đặt A $=-1$, B $=1$ với \sig{ALUCtrl}=8 rồi 9: SLT cho 1, SLTU cho 0. Đối chiếu \texttt{tests/alu.vec}.\end{thinghiem}
\begin{kiemtra}Vì sao kết quả SLT phải mở rộng 0? Bộ phận nào giúp \texttt{sra} giữ dấu?\end{kiemtra}

\subsection{Khối \texorpdfstring{\texttt{BranchUnit}}{BranchUnit}}
\begin{itemize}
  \item \textbf{Mục đích:} quyết định branch có được rẽ hay không.
  \item \textbf{Vào / ra:} vào \sig{rs1_data}, \sig{rs2_data}, \sig{IsBranch}, \sig{funct3}; ra \sig{branch_taken}.
  \item \textbf{Bên trong:} một bộ so sánh không dấu và một có dấu 32 bit (cho $=$, $<_u$, $<_s$); 3 cổng NOT (cho $\ne$, $\ge_u$, $\ge_s$); multiplexer 8-ra-1 một bit chọn theo \sig{funct3} (000 beq, 001 bne, 100 blt, 101 bge, 110 bltu, 111 bgeu; 010 và 011 cho 0); cổng AND với \sig{IsBranch}.
  \item \textbf{Vì sao:} sáu điều kiện là ba phép so sánh và phủ định của chúng; AND với \sig{IsBranch} đảm bảo lệnh khác không bao giờ rẽ.
\end{itemize}
\begin{kiemtra}\texttt{blt} và \texttt{bltu} khác nhau trong mạch ở chỗ nào?\end{kiemtra}

\textbf{Bài kiểm tra cuối giai đoạn:} tính tay \texttt{sub x13,x12,x10} qua \sig{RegisterFile} $\rightarrow$ \sig{ALUControl} $\rightarrow$ \sig{ALU}, rồi kiểm chứng bằng demo (tick 38).

% =====================================================================
\section{Giai đoạn 4 --- Bộ nhớ và CPU single-cycle (buổi 6--7)}
% =====================================================================

\subsection{Khối \texorpdfstring{\texttt{Memory}}{Memory} cùng \texorpdfstring{\texttt{LoadExtract}}{LoadExtract} và \texorpdfstring{\texttt{StoreMerge}}{StoreMerge}}
\begin{itemize}
  \item \textbf{Mục đích:} bộ nhớ dữ liệu với truy cập byte, nửa từ và từ (\texttt{lb lh lw lbu lhu sb sh sw}).
  \item \textbf{Vào / ra:} vào \sig{Address}, \sig{write_data}, \sig{MemWrite}, \sig{CLK}, \sig{funct3}; ra \sig{read_data}.
  \item \textbf{Bên trong:} một RAM 256$\times$32 với \emph{đọc bất đồng bộ} và một chân cho phép ghi (line enables); một splitter địa chỉ (độ lệch byte \sig{Address[1:0]}, chỉ số từ \sig{Address[9:2]}); \sig{LoadExtract}; \sig{StoreMerge}.
  \item \textbf{Vì sao đọc bất đồng bộ:} lệnh load phải có dữ liệu ngay trong chu kỳ. Với đọc đồng bộ mặc định, load trả về dữ liệu không xác định (một trong các lỗi đã gặp).
\end{itemize}

\textbf{\sig{LoadExtract}:} splitter (4 byte, 2 nửa từ), multiplexer 4-ra-1$\times$8 và 2-ra-1$\times$16 chọn byte / nửa từ, 4 bộ mở rộng bit (dấu/0, 8 và 16 $\rightarrow$ 32), multiplexer 8-ra-1$\times$32 cuối cùng chọn theo \sig{funct3}. \texttt{lbu}/\texttt{lhu} mở rộng 0, \texttt{lb}/\texttt{lh} mở rộng dấu.

\textbf{\sig{StoreMerge}:} RAM chỉ ghi nguyên từ, nên store byte là \textbf{đọc--sửa--ghi}: đọc từ cũ, thay một làn byte, ghi cả từ lại.
Một ROM 16$\times$4 biến $\{\mathrm{size},\mathrm{addr}[1:0]\}$ thành bốn tín hiệu cho phép byte (\texttt{1 2 4 8} cho byte, \texttt{3 3 c c} cho nửa từ, \texttt{f} cho từ); bốn multiplexer 2-ra-1 8 bit chọn, theo từng làn, byte mới hay byte cũ; thêm 3 multiplexer byte chọn nguồn byte cho từng làn.

\begin{figure}[htbp]
\centering
\begin{tikzpicture}[node distance=8mm]
  \node[mem] (ram) {RAM\\256$\times$32};
  \node[blk,right=16mm of ram,yshift=8mm] (le) {\sig{LoadExtract}\\chọn byte / nửa từ\\mở rộng dấu / 0};
  \node[lab,right=8mm of le] (rd) {\sig{read_data}};
  \node[blk,below=14mm of ram] (sm) {\sig{StoreMerge}\\(từ cũ + dữ liệu mới)};
  \node[lab,left=10mm of sm] (wd) {\sig{write_data}};
  \node[lab,left=10mm of ram,yshift=6mm] (ad) {\sig{Address}};
  \draw[ar] (ad) -- (ram);
  \draw[ar] (ram) -- node[lab,above]{từ cũ} (le);
  \draw[bus] (le) -- (rd);
  \draw[ar] (ram.south) -- node[lab,right]{từ cũ} (sm.north);
  \draw[ar] (wd) -- (sm);
  \draw[ar] (sm.east) -- ++(10mm,0) |- node[lab,pos=.2,right]{từ mới (ghi)} ($(ram.east)+(0,-4mm)$);
\end{tikzpicture}
\caption{\sig{Memory}: đọc đi qua \sig{LoadExtract}; ghi đi qua \sig{StoreMerge} (đọc--sửa--ghi).}
\end{figure}

\begin{thinghiem}Trong \texttt{tests/loadextract.vec} và \texttt{tests/storemerge.vec}, tìm các dòng \texttt{sb} ở độ lệch 2 và \texttt{lb} của byte có bit cao bằng 1.\end{thinghiem}
\begin{kiemtra}Vì sao \sig{StoreMerge} đọc được từ cũ? Với từ cũ \texttt{0x11223344} và $x1=\mathtt{0xAA}$, lệnh \texttt{sb x1,1(x0)} ghi gì vào từ 0? (Đáp án: \texttt{0x1122AA44}.)\end{kiemtra}

\subsection{CPU single-cycle \texorpdfstring{\texttt{main}}{main}}
\begin{itemize}
  \item \textbf{Bên trong:} thanh ghi PC (xóa bằng \sig{RESET}); ROM lệnh 256$\times$32 đánh địa chỉ bằng \sig{PC[9:2]}; 2 bộ cộng (\sig{PC+4}, \sig{PC+imm}); cổng AND 32 bit với mặt nạ \texttt{0xFFFFFFFE} (cho JALR); một cổng OR (\sig{branch_taken} hoặc \sig{IsJAL}); 5 multiplexer (toán hạng A theo \sig{ASel}, toán hạng B theo \sig{ALUSrc}, write-back theo \sig{WBSel}, hai cái cho PC kế tiếp); 8 sub-circuit.
  \item \textbf{PC kế tiếp:} $\mathtt{next\_pc}=\mathtt{IsJALR}\ ?\ \mathtt{jalr\_target}:(\mathtt{branch\_taken}\lor\mathtt{IsJAL})\ ?\ \mathtt{PC{+}imm}:\mathtt{PC{+}4}$.
  \item \textbf{Vì sao thiết kế thế:} \texttt{lui}, \texttt{auipc}, \texttt{jal}, \texttt{jalr} dùng lại đường ALU và write-back nhờ \sig{ASel} và \sig{WBSel}, không cần phần cứng mới.
\end{itemize}

\begin{figure}[p]
\centering
\begin{tikzpicture}[x=10mm,y=10mm,scale=0.86,transform shape]
  % IF
  \node[blk,minimum width=10mm] (pc) at (0,0) {PC};
  \node[blk,minimum width=10mm] (add4) at (0,-2) {+4};
  \node[mem] (rom) at (2.3,0) {ROM lệnh};
  % ID
  \node[blk] (dec) at (4.6,0) {Giải mã};
  \node[ctl] (ctl) at (4.6,1.6) {\sig{MainControl}};
  \node[blk] (rf) at (7,0) {\sig{RegisterFile}};
  \node[blk] (imm) at (7,-1.8) {\sig{ImmGen}};
  % EX
  \node[mux] (ma) at (9.1,0.4) {A};
  \node[mux] (mb) at (9.1,-1.0) {B};
  \node[blk,minimum height=22mm,minimum width=12mm] (alu) at (10.8,-0.3) {\sig{ALU}};
  \node[ctl] (actl) at (10.8,1.6) {\sig{ALUControl}};
  \node[blk] (br) at (10.8,-2.4) {\sig{BranchUnit}};
  \node[blk] (addt) at (10.8,-3.5) {$\mathtt{PC}+\mathtt{imm}$};
  % MEM / WB
  \node[mem] (mem) at (14.0,-0.3) {\sig{Memory}};
  \node[mux,minimum width=16mm] (wb) at (15.8,-0.3) {WB};
  % IF
  \draw[ar] (pc) -- (rom);
  \draw[ar] (pc) -- (add4);
  \draw[ar] (add4.west) -- ++(-0.6,0) |- node[lab,pos=.25,left]{\sig{next_pc}} (pc.west);
  % ID
  \draw[bus] (rom) -- node[lab,above]{lệnh} (dec);
  \draw[ar] (dec) -- (ctl);
  \draw[ar] (dec) -- (rf);
  \draw[ar] (dec.south) |- (imm.west);
  % EX
  \draw[bus] (rf.east) -- ++(0.4,0) |- (ma.west);
  \draw[bus] ($(rf.east)+(0.4,0)$) |- ($(mb.west)+(0,0.15)$);
  \draw[bus] (imm.east) -- ++(0.9,0) |- ($(mb.west)+(0,-0.15)$);
  \draw[bus] (ma.east) -- (ma.east -| alu.west);
  \draw[bus] (mb.east) -- (mb.east -| alu.west);
  \draw[cs] (actl) -- (alu);
  \draw[ar] (br.east) -- ++(0.8,0) node[lab,right]{\sig{branch_taken}};
  \draw[ar] (addt.east) -- ++(0.8,0) node[lab,right]{\sig{pc_target}};
  % MEM / WB
  \draw[bus] (alu) -- node[lab,above,pos=.45]{\sig{alu_result}} (mem);
  \draw[bus] (mem) -- (wb);
  \draw[bus] ($(alu.east)+(0.3,0)$) |- ++(0,-1.3) -| ($(wb.south)$);
  \draw[bus] (wb.east) -- ++(0.3,0) |- (16.1,2.4) -| node[lab,pos=.3,above]{\sig{wb_data} ghi về \sig{RegisterFile}} (rf.north);
  \draw[cs] (ctl.west) -- ++(-1.0,0) node[lab,left,align=right]{tín hiệu\\điều khiển};
\end{tikzpicture}
\caption{Sơ đồ khối CPU single-cycle (rút gọn): nạp lệnh $\rightarrow$ giải mã $\rightarrow$ thực thi $\rightarrow$ bộ nhớ $\rightarrow$ ghi lại. Nét cam đứt: tín hiệu điều khiển.}
\label{fig:main}
\end{figure}

\textbf{Một lệnh đi qua năm tầng (ví dụ cho \texttt{add x10,x10,x11}, tick 3):}
\begin{enumerate}
  \item \textbf{IF:} ROM tại PC = \texttt{0x08} xuất \texttt{00b50533}; \sig{pc_plus4} = \texttt{0x0c}.
  \item \textbf{ID:} cắt $\mathtt{opcode}=\mathtt{0x33}$, $\mathtt{rd}=10$, $\mathtt{rs1}=10$, $\mathtt{rs2}=11$. \sig{MainControl}: $\mathtt{RegWrite}=1$, $\mathtt{ALUSrc}=0$, $\mathtt{ALUOp}=10$. Register file xuất $x10=0$ và $x11=\mathtt{0x0a}$.
  \item \textbf{EX:} \sig{ALUControl} cho $\mathtt{ALUCtrl}=0$ (ADD); ALU tính $0+\mathtt{0x0a}=\mathtt{0x0a}$.
  \item \textbf{MEM:} không có gì.
  \item \textbf{WB:} $\mathtt{wb\_data}=\mathtt{0x0a}$ được ghi vào $x10$ ở cạnh xung nhịp; PC $=$ PC$+4=\mathtt{0x0c}$.
\end{enumerate}

\begin{thinghiem}Dùng các probe ở cuối mạch và bảng ở phần~\ref{sec:demo}: theo dõi một lệnh mỗi lớp (add, addi, bnez, sw, lw, jal, ret) qua năm tầng với giá trị thật.\end{thinghiem}
\begin{kiemtra}Vẽ lại \sig{main} từ trí nhớ thành các hộp IF, ID, EX, MEM, WB và đánh dấu ba tín hiệu chọn multiplexer lấy từ đâu. Lệnh nào dùng bộ cộng riêng \sig{PC+imm}? Đích của JALR tính thế nào?\end{kiemtra}

% =====================================================================
\section{Giai đoạn 5 --- Pipeline (buổi 8--10)}
% =====================================================================

\subsection{Các tầng, thanh ghi pipeline, thời điểm ghi}
Năm tầng làm việc trên năm lệnh cùng lúc. Bốn \textbf{thanh ghi pipeline} (IF/ID, ID/EX, EX/MEM, MEM/WB) mang những gì tầng sau cần.

\begin{figure}[htbp]
\centering
\begin{tikzpicture}[node distance=3mm]
  \node[blk,minimum width=17mm] (sif) {IF\\nạp lệnh};
  \node[reg,right=of sif] (r1) {};
  \node[blk,minimum width=17mm,right=of r1] (id) {ID\\giải mã};
  \node[reg,right=of id] (r2) {};
  \node[blk,minimum width=17mm,right=of r2] (ex) {EX\\ALU};
  \node[reg,right=of ex] (r3) {};
  \node[blk,minimum width=17mm,right=of r3] (me) {MEM\\RAM};
  \node[reg,right=of me] (r4) {};
  \node[blk,minimum width=17mm,right=of r4] (wb) {WB\\ghi lại};
  \foreach \a/\b in {sif/r1,r1/id,id/r2,r2/ex,ex/r3,r3/me,me/r4,r4/wb} \draw[ar] (\a) -- (\b);
  \node[lab,below=2mm of r1] {IF/ID};
  \node[lab,below=2mm of r2] {ID/EX};
  \node[lab,below=2mm of r3] {EX/MEM};
  \node[lab,below=2mm of r4] {MEM/WB};
  \node[ctl,above=12mm of ex] (fu) {\sig{ForwardingUnit}};
  \node[ctl,above=12mm of id] (hu) {\sig{HazardUnit}};
  \draw[cs] (fu.south) -- (ex.north);
  \draw[cs] (hu.south) -- (id.north);
  \draw[cs] (r3.north) |- (fu.east);
  \draw[cs] (r4.north) |- ($(fu.east)+(0,2mm)$);
  \draw[cs] (hu.west) -| node[lab,pos=.75,left,align=right]{\sig{run}:\\giữ PC,\\IF/ID} (sif.north);
  \draw[cs] (hu.south west) -- (r1.north);
\end{tikzpicture}
\caption{Pipeline 5 tầng, bốn thanh ghi pipeline, hai khối xử lý hazard.}
\label{fig:pipe}
\end{figure}

\begin{table}[htbp]
\centering
\small
\begin{tabular}{@{}lp{10cm}@{}}
\toprule
\textbf{Thanh ghi} & \textbf{Chứa}\\
\midrule
IF/ID & \sig{pc_id}, \sig{instr_id}\\
ID/EX & \sig{pc_ex}, \sig{rs1_data_ex}, \sig{rs2_data_ex}, \sig{imm_ex}, \sig{rs1_ex}, \sig{rs2_ex}, \sig{rd_ex}, \sig{funct3_ex}, \sig{funct7_ex}, từ điều khiển \sig{ctrl_ex} (13 bit: \sig{RegWrite ALUSrc MemWrite MemRead IsBranch ALUOp ASel WBSel IsJAL IsJALR})\\
EX/MEM & \sig{alu_result_mem}, \sig{store_data_mem}, \sig{pc_plus4_mem}, \sig{rd_mem}, \sig{funct3_mem}, điều khiển \sig{RegWrite MemWrite MemRead WBSel}\\
MEM/WB & \sig{alu_result_wb}, \sig{mem_rdata_wb}, \sig{pc_plus4_wb}, \sig{rd_wb}, điều khiển \sig{RegWrite WBSel}\\
\bottomrule
\end{tabular}
\caption{Nội dung bốn thanh ghi pipeline (hậu tố \sig{_id}, \sig{_ex}, \sig{_mem}, \sig{_wb} là tầng hiện tại).}
\end{table}

\textbf{Bên trong \sig{Pipeline}:} 24 thanh ghi (PC cùng 23 thanh ghi pipeline có nhãn như \sig{pc_id}, \sig{alu_result_mem}); 8 splitter (các bit PC, đóng gói / tách từ điều khiển); 10 sub-circuit; 10 multiplexer (xem phần~\ref{sec:circ}); 3 bộ cộng (\sig{PC+4} ở IF, \sig{pc_ex+imm} và \sig{pc_ex+4} ở EX); hai cổng NOT: một tạo \sig{CLK_n}, một tạo $\mathtt{run}=\neg\mathtt{stall}$; hai cổng OR (\sig{pc_src_ex} và \sig{bubble}).

\textbf{Từ điều khiển:} tín hiệu điều khiển được sinh một lần ở ID rồi đóng gói; mỗi tầng chỉ tách phần mình cần.

\textbf{Thời điểm ghi:} register file ghi ở \textbf{sườn xuống} (\sig{CLK_n}), nên lệnh ở WB ghi trong nửa đầu chu kỳ và lệnh ở ID đọc giá trị mới ở nửa sau: không cần forwarding cho khoảng cách 3 lệnh.

\begin{thinghiem}Mở \sig{Pipeline}, tìm bốn nhóm thanh ghi theo nhãn; với \texttt{add x2,x1,x1} hãy ghi lại nội dung từng nhóm sau 1, 2, 3, 4 tick.\end{thinghiem}
\begin{kiemtra}Vì sao PC và IF/ID có chân cho phép ghi \sig{run} còn các thanh ghi khác thì không? Vì sao các thanh ghi pipeline bằng 0 lúc đầu và điều đó vô hại?\end{kiemtra}

\subsection{Khối \texorpdfstring{\texttt{ForwardingUnit}}{ForwardingUnit}}
\begin{itemize}
  \item \textbf{Mục đích:} đưa cho EX giá trị mới nhất của \sig{rs1}/\sig{rs2} khi lệnh trước chưa kịp ghi.
  \item \textbf{Vào / ra:} vào \sig{rs1_ex}, \sig{rs2_ex}, \sig{rd_mem}, \sig{RegWrite_mem}, \sig{rd_wb}, \sig{RegWrite_wb}; ra \sig{ForwardA}, \sig{ForwardB} (2 bit).
  \item \textbf{Bên trong:} 6 bộ so sánh 5 bit, các cổng AND/NOT, 2 splitter tạo tín hiệu chọn 2 bit.
\end{itemize}
\textbf{Quy tắc:}
\[
\mathtt{ForwardA}=
\begin{cases}
10 & \text{nếu } \mathtt{RegWrite\_mem}\ \wedge\ \mathtt{rd\_mem}\neq0\ \wedge\ \mathtt{rd\_mem}=\mathtt{rs1\_ex}\ \ (\text{từ EX/MEM})\\
01 & \text{ngược lại nếu điều tương tự đúng với WB}\ \ (\text{từ MEM/WB})\\
00 & \text{còn lại}\ \ (\text{register file})
\end{cases}
\]
EX/MEM được ưu tiên vì mới hơn; điều kiện $\mathtt{rd}\neq0$ ngăn forward \sig{x0}.
\textbf{Hai chi tiết:} giá trị lấy từ EX/MEM là kết quả ALU \emph{trừ JAL/JALR thì là PC+4}; toán hạng B sau forwarding cũng là dữ liệu của lệnh store.

\begin{figure}[htbp]
\centering
\begin{tikzpicture}[x=10mm,y=10mm]
  \node[mux,minimum width=18mm,minimum height=8mm] (ma) at (4.5,1.1) {A};
  \node[mux,minimum width=18mm,minimum height=8mm] (mb) at (4.5,-1.1) {B};
  \node[blk,minimum height=30mm,minimum width=12mm] (alu) at (6.8,0) {\sig{ALU}};
  \node[ctl] (fa) at (4.5,2.7) {\sig{ForwardA}};
  \node[ctl] (fb) at (4.5,-2.7) {\sig{ForwardB}};
  \foreach \m in {ma,mb} {
    \foreach \dy/\t in {0.45/{00: register file (đọc ở ID)},0/{10: EX/MEM (ALU hoặc PC+4)},-0.45/{01: MEM/WB (dữ liệu ghi lại)}} {
      \draw[ar] ($(\m.west)+(-1.2,\dy)$) node[lab,left]{\t} -- ($(\m.west)+(0,\dy)$);
    }
  }
  \draw[bus] (ma.east) -- (ma.east -| alu.west);
  \draw[bus] (mb.east) -- (mb.east -| alu.west);
  \draw[cs] (fa) -- (ma);
  \draw[cs] (fb) -- (mb);
  \draw[ar] ($(mb.east)+(0.3,0)$) |- ++(0,-1.0) -- ++(3.0,0) node[lab,right]{dữ liệu store};
\end{tikzpicture}
\caption{Forwarding: toán hạng của ALU lấy từ register file (00), MEM/WB (01) hoặc EX/MEM (10).}
\end{figure}

\subsection{Khối \texorpdfstring{\texttt{HazardUnit}}{HazardUnit}}
\begin{itemize}
  \item \textbf{Mục đích:} phát hiện lệnh load theo sau ngay bởi lệnh dùng kết quả của nó.
  \item \textbf{Vào / ra:} vào \sig{MemRead_ex}, \sig{rd_ex}, \sig{rs1_id}, \sig{rs2_id}; ra \sig{stall}.
  \item \textbf{Bên trong:} 3 bộ so sánh 5 bit, cổng OR, NOT và AND 3 ngõ vào.
  \item \textbf{Quy tắc:} $\mathtt{stall}=\mathtt{MemRead\_ex}\wedge\mathtt{rd\_ex}\ne0\wedge(\mathtt{rd\_ex}=\mathtt{rs1\_id}\lor\mathtt{rd\_ex}=\mathtt{rs2\_id})$.
  \item \textbf{Tác dụng:} $\mathtt{run}=\neg\mathtt{stall}$ làm PC và IF/ID đứng yên, còn từ điều khiển vào ID/EX bị thay bằng 0 (một bubble).
\end{itemize}

\begin{kiemtra}Vì sao load không forward được như kết quả ALU? Vì sao chỉ stall một chu kỳ mà không phải hai? Bộ so sánh nào bắt điều kiện $\mathtt{rd\_ex}=\mathtt{rs2\_id}$?\end{kiemtra}

\subsection{Hazard điều khiển và hiệu năng}
Branch, JAL, JALR được quyết định ở \textbf{EX}: $\mathtt{pc\_src\_ex}=\mathtt{branch\_taken}\lor\mathtt{IsJAL}\lor\mathtt{IsJALR}$. PC kế tiếp thành địa chỉ đích, lệnh đang nạp bị thay bằng 0 ($\mathtt{instr\_in}=0$) và từ điều khiển vào ID/EX bị xóa: hai lệnh sau lệnh rẽ nhánh bị hủy. Chi phí: 2 chu kỳ cho mỗi branch/lệnh nhảy được rẽ; branch không rẽ không tốn gì (dự đoán tĩnh ``không rẽ'').

\begin{center}
\begin{tabular}{@{}lc@{}}
\toprule
\textbf{Tình huống} & \textbf{Chu kỳ thêm}\\
\midrule
lệnh phụ thuộc ngay sau lệnh ALU (forwarding) & 0\\
dùng giá trị vừa load ở lệnh kế tiếp (load-use) & 1\\
load, một lệnh không liên quan, rồi mới dùng & 0\\
branch không rẽ & 0\\
branch được rẽ / \texttt{jal} / \texttt{jalr} & 2\\
\bottomrule
\end{tabular}
\end{center}

Tính theo chu kỳ, pipeline ở đây không nhanh hơn single-cycle, vì mô phỏng Logisim không có độ trễ cổng. Lợi ích thật là \textbf{chu kỳ xung nhịp ngắn hơn} (đường dài nhất còn một tầng thay vì năm). Cách cải thiện: quyết định JAL ở ID, bộ dự đoán nhánh, bộ đệm đích nhánh.

\textbf{Giải thích 64 chu kỳ của \texttt{demo.s}:} 39 lệnh $+$ 11 lần chuyển hướng $\times2=22$ (9 \texttt{bnez} được rẽ, \texttt{jal}, \texttt{ret}) $=61$; 3 chu kỳ còn lại là thời gian làm đầy pipeline.

\subsubsection*{Chương trình \texttt{hazards\_demo.s} trên pipeline}
\begin{center}
\small
\begin{tabular}{@{}lll@{}}
\toprule
\textbf{Địa chỉ} & \textbf{Lệnh} & \textbf{Mục đích}\\
\midrule
0x00 & \texttt{li x1,100} & x1 = 100\\
0x04 & \texttt{sw x1,0(x0)} & cần x1 ngay: \textbf{forwarding}\\
0x08 & \texttt{lw x2,0(x0)} & load\\
0x0c & \texttt{addi x3,x2,1} & dùng x2 ngay: \textbf{stall load-use}\\
0x10 & \texttt{add x4,x3,x2} & x3 được forward\\
0x14 & \texttt{beq x4,x4,skip} & luôn rẽ: \textbf{flush}\\
0x18 & \texttt{addi x5,x0,1} & không được có hiệu lực\\
0x1c & \texttt{addi x6,x0,2} & không được có hiệu lực\\
0x20 & \texttt{skip: addi x7,x0,7} & lệnh đầu tiên ở đích\\
0x24 & \texttt{halt: j halt} & kết thúc\\
\bottomrule
\end{tabular}
\end{center}
Kết quả: $x1=\mathtt{0x64}$, $x2=\mathtt{0x64}$, $x3=\mathtt{0x65}$, $x4=\mathtt{0xc9}$, $x7=7$ và $x5=x6=0$.

\begin{figure}[p]
\centering
\begin{tikzpicture}[x=11mm,y=-6.5mm]
  % cột thời gian
  \foreach \cl in {0,...,13} \node[lab] at (\cl+0.5,-0.2) {\cl};
  \node[lab,anchor=east] at (-0.1,-0.2) {sau tick};
  % hàng: lệnh
  \node[lab,anchor=east] at (-0.1,1.2) {\texttt{li x1}};
  \node[lab,anchor=east] at (-0.1,2.2) {\texttt{sw x1}};
  \node[lab,anchor=east] at (-0.1,3.2) {\texttt{lw x2}};
  \node[lab,anchor=east] at (-0.1,4.2) {\texttt{addi x3}};
  \node[lab,anchor=east] at (-0.1,5.2) {\texttt{add x4}};
  \node[lab,anchor=east] at (-0.1,6.2) {\texttt{beq}};
  \node[lab,anchor=east] at (-0.1,7.2) {\texttt{addi x5}};
  \node[lab,anchor=east] at (-0.1,8.2) {\texttt{addi x6}};
  \node[lab,anchor=east] at (-0.1,9.2) {\texttt{addi x7}};
  % các ô: (hàng, cột, nhãn)
  \foreach \rr/\cl/\lbl in {
    1/0/IF,1/1/ID,1/2/EX,1/3/ME,1/4/WB,
    2/1/IF,2/2/ID,2/3/EX,2/4/ME,2/5/WB,
    3/2/IF,3/3/ID,3/4/EX,3/5/ME,3/6/WB,
    4/3/IF,4/4/ID,4/5/ID,4/6/EX,4/7/ME,4/8/WB,
    5/4/IF,5/5/IF,5/6/ID,5/7/EX,5/8/ME,5/9/WB,
    6/6/IF,6/7/ID,6/8/EX,6/9/ME,6/10/WB,
    7/7/IF,7/8/ID,
    8/8/IF,
    9/9/IF,9/10/ID,9/11/EX,9/12/ME,9/13/WB} {
    \node[stg,fill=blue!8,minimum width=10.5mm] at (\cl+0.5,\rr+0.2) {\tiny\lbl};}
  % stall và flush
  \node[stg,fill=orange!35,minimum width=10.5mm] at (5.5,4.2) {\tiny ID};
  \node[stg,fill=red!25,minimum width=10.5mm] at (9.5,7.2) {\tiny X};
  \node[stg,fill=red!25,minimum width=10.5mm] at (9.5,8.2) {\tiny X};
  % chú thích
  \node[stg,fill=orange!35,minimum width=8mm] (l1) at (0.5,10.6) {\tiny ID};
  \node[lab,anchor=west] at (l1.east) {\ stall load-use (giữ ở ID)};
  \node[stg,fill=red!25,minimum width=8mm] (l2) at (5.5,10.6) {\tiny X};
  \node[lab,anchor=west] at (l2.east) {\ bị flush (thành bubble)};
\end{tikzpicture}
\caption{Pipeline của \texttt{hazards\_demo.s}. Cột là trạng thái sau mỗi tick. Cam: lệnh \texttt{addi x3} bị giữ ở ID một chu kỳ (stall load-use). Đỏ (X): hai lệnh sau \texttt{beq} bị flush thành bubble.}
\label{fig:hazard}
\end{figure}

\subsubsection*{Trạng thái các tầng sau mỗi tick (đo trên project)}
\begin{center}
\small
\setlength{\tabcolsep}{3pt}
\begin{tabular}{@{}rllllcccc@{}}
\toprule
\textbf{Tick} & \textbf{IF (PC)} & \textbf{ID} & \textbf{EX} & \textbf{MEM} & \sig{stall} & \sig{pc_src_ex} & \sig{FwdA} & \sig{FwdB}\\
\midrule
1 & sw (04) & li & -- & -- & 0 & 0 & 00 & 00\\
2 & lw (08) & sw & li & -- & 0 & 0 & 00 & 00\\
3 & addi (0c) & lw & sw & li & 0 & 0 & 00 & \textbf{10}\\
4 & add (10) & addi & lw & sw & \textbf{1} & 0 & 00 & 00\\
5 & add (\textbf{10}, giữ) & addi & bubble & lw & 0 & 0 & 10 & 00\\
6 & beq (14) & add & addi & bubble & 0 & 0 & \textbf{01} & 00\\
7 & addi x5 (18) & beq & add & addi & 0 & 0 & \textbf{10} & 00\\
8 & addi x6 (1c) & addi x5 & \textbf{beq} & add & 0 & \textbf{1} & 10 & 10\\
9 & skip (\textbf{20}) & bubble & bubble & beq & 0 & 0 & 00 & 00\\
10 & halt (24) & addi x7 & bubble & bubble & 0 & 0 & 00 & 00\\
11 & (28) & halt & addi x7 & bubble & 0 & 0 & 00 & 00\\
12 & (2c) & bubble & halt & addi x7 & 0 & \textbf{1} & 00 & 00\\
\bottomrule
\end{tabular}
\end{center}

\begin{thinghiem}Nạp \texttt{programs/hazards\_demo.hex} vào \sig{Pipeline}, bấm 12 tick và tự điền bảng trên bằng các probe \sig{stall}, \sig{pc_src_ex}, \sig{ForwardA}, \sig{ForwardB} trước khi đọc đáp án.\end{thinghiem}
\begin{kiemtra}Vì sao $x5=x6=0$ là bằng chứng flush hoạt động? Nếu không có flush thì sao? Vì sao \sig{ForwardA}=10 ở tick 5 vô hại?\end{kiemtra}

\textbf{Bài kiểm tra cuối giai đoạn:} vẽ giản đồ pipeline (IF ID EX ME WB theo từng chu kỳ) của \texttt{hazards\_demo.s} trên giấy, rồi so với Hình~\ref{fig:hazard}.

% =====================================================================
\section{Giai đoạn 6 --- Chương trình, kiểm thử, công cụ (buổi 11--12)}
% =====================================================================

\subsection{Chương trình và công cụ}
\begin{itemize}
  \item \textbf{Assembler} (\texttt{tools/rv.py asm}): nhãn, lệnh giả (\texttt{li}, \texttt{mv}, \texttt{bnez}, \texttt{ret}, \texttt{j}), macro kiểm tra; tạo ảnh ROM dạng \texttt{v2.0 raw}.
  \item \textbf{Bộ mô phỏng tham chiếu} (\texttt{tools/rv.py sim}): viết độc lập với mạch; là ``chuẩn'' của mọi bài test.
  \item \textbf{Bộ chạy không giao diện} (\texttt{tools/java/Run.java}): nạp ROM, đảo xung nhịp, đọc PC, thanh ghi và RAM của mạch thật.
  \item \textbf{Bộ sinh mạch} (\texttt{tools/build}): sinh circuit từ netlist và kiểm tra từng net.
\end{itemize}

\subsection{Ba mức kiểm thử}
\begin{enumerate}
  \item \textbf{Test module} (\texttt{tests/run\_tests.ps1}): 106 hàng test-vector trên 9 module, do chính Logisim chạy.
  \item \textbf{Chương trình} (\texttt{tests/run\_cpu\_tests.py}): 7 chương trình trên cả hai CPU; sau khi chạy, cả 32 thanh ghi và mọi từ dữ liệu khác 0 phải bằng kết quả của bộ mô phỏng tham chiếu.
  \item \textbf{Chương trình ngẫu nhiên} (\texttt{tests/fuzz.py}): chương trình ngẫu nhiên 60--80 lệnh chỉ dùng 7 thanh ghi (nên hầu hết lệnh đều phụ thuộc lệnh gần đó).
\end{enumerate}

\begin{figure}[htbp]
\centering
\begin{tikzpicture}[node distance=10mm]
  \node[blk] (src) {chương trình\\\texttt{.s}};
  \node[blk,right=of src] (asm) {assembler\\\texttt{rv.py asm}};
  \node[mem,right=of asm] (hex) {ảnh ROM\\\texttt{.hex}};
  \node[blk,above right=4mm and 14mm of hex] (c1) {mạch \sig{main}};
  \node[blk,below right=4mm and 14mm of hex] (c2) {mạch \sig{Pipeline}};
  \node[blk,right=of src,yshift=-18mm,xshift=0mm] (sim) {mô phỏng\\tham chiếu};
  \node[ctl,right=30mm of hex] (cmp) {so sánh 32\\thanh ghi + RAM\\$\Rightarrow$ \textbf{OK} / \textbf{DIFFERENT}};
  \draw[ar] (src) -- (asm); \draw[ar] (asm) -- (hex);
  \draw[ar] (hex) -- (c1); \draw[ar] (hex) -- (c2);
  \draw[ar] (c1) -- (cmp); \draw[ar] (c2) -- (cmp);
  \draw[ar] (src.south) |- (sim.west);
  \draw[ar] (sim.east) -| (cmp.south);
  \end{tikzpicture}
\caption{Cách chứng minh: cùng một chương trình chạy trên hai mạch và trên bộ mô phỏng độc lập, rồi so sánh.}
\end{figure}

\begin{thinghiem}Viết \texttt{test1.s} (giai thừa của 5), chạy \texttt{python tools/live\_proof.py test1.s}; sau đó sửa một lệnh và \emph{dự đoán} kết quả trước khi chạy.\end{thinghiem}
\begin{kiemtra}Vì sao bộ mô phỏng tham chiếu phải độc lập với mạch? Bộ fuzz tìm ra điều gì mà test viết tay không tìm?\end{kiemtra}

\subsection{Lỗi đã tìm thấy và bài học}
\begin{center}
\small
\begin{tabular}{@{}lp{5.2cm}p{4.5cm}@{}}
\toprule
\textbf{Ở đâu} & \textbf{Lỗi} & \textbf{Cách sửa}\\
\midrule
ALU & \texttt{SLTU} trả \texttt{FFFFFFFF} thay vì 1 (mở rộng dấu) & mở rộng 0\\
ALUControl & \texttt{SLL} kiểu R chọn mã của \texttt{SRL} & sửa hằng số\\
MainControl & JALR so với opcode của JAL; LUI chọn hằng số kiểu B & viết lại phần giải mã\\
ImmGen & kiểu U thiếu $\ll12$; bit 0 kiểu J bằng 1 & đã sửa\\
Memory & load trả dữ liệu không xác định (RAM đọc đồng bộ) & đọc bất đồng bộ\\
Datapath & JAL, JALR, LUI, AUIPC chưa có đường dữ liệu & thêm \sig{ASel}, \sig{WBSel}, logic PC\\
\bottomrule
\end{tabular}
\end{center}
\textbf{Cạm bẫy của Logisim:} dây chỉ nối ở \emph{đầu mút} (không nối ở chỗ giao chữ T); mọi thứ phải nằm trên lưới 10 px; \sig{Constant} mặc định bằng 1; RAM mặc định đọc đồng bộ.

% =====================================================================
\section{Đọc sâu file mạch \texttt{RiscVComputer.circ}}
\label{sec:circ}
% =====================================================================
Phần này lấy trực tiếp từ file \texttt{.circ} (Logisim Evolution 4.1.0): số linh kiện thật, cách các dây và tunnel nối với nhau, và những chi tiết chỉ thấy khi lần từng net. Nên đọc sau giai đoạn 5.

\subsection{Số linh kiện trong từng circuit}
\begin{center}
\small
\begin{tabularx}{\linewidth}{@{}l>{\raggedright\arraybackslash}X@{}}
\toprule
\textbf{Circuit} & \textbf{Linh kiện (không kể Pin, Tunnel, Text, Probe)}\\
\midrule
\sig{InstructionDecode} & 1 splitter (32 $\rightarrow$ 6 trường)\\
\sig{MainControl} & 9 bộ so sánh 7 bit + 9 hằng số opcode; 11 cổng OR (8 cổng 3 ngõ vào, 3 cổng 2 ngõ vào); 4 splitter ghép bus\\
\sig{ImmGen} & 9 splitter; 4 bộ mở rộng bit; 3 hằng số (hai 1 bit, một 12 bit); 1 MUX 8:1\\
\sig{RegisterFile} & 32 thanh ghi; 1 DEMUX 1:32; 2 MUX 32:1; 1 hằng số 0\\
\sig{ALUControl} & 19 hằng số 4 bit; 5 MUX (hai 8:1, một 4:1, hai 2:1); 1 splitter tách \sig{funct7[5]}\\
\sig{ALU} & bộ cộng, bộ trừ, AND/OR/XOR 32 bit, 3 bộ dịch, 2 bộ so sánh, 2 bộ mở rộng 0, 1 MUX 16:1, 1 splitter lấy \sig{B[4:0]}\\
\sig{BranchUnit} & 2 bộ so sánh 32 bit; 3 NOT; 1 MUX 8:1 một bit; 1 AND; 1 hằng số 0\\
\sig{Memory} & 1 RAM 256$\times$32 (\sig{asyncread=true}); 1 splitter địa chỉ; \sig{LoadExtract}; \sig{StoreMerge}\\
\sig{LoadExtract} & 3 splitter; 3 MUX (4:1$\times$8, 2:1$\times$16, 8:1$\times$32); 4 bộ mở rộng bit\\
\sig{StoreMerge} & 6 splitter; 1 ROM 16$\times$4; 7 MUX (ba 4:1$\times$8 chọn nguồn byte, bốn 2:1$\times$8 chọn byte mới/cũ)\\
\sig{ForwardingUnit} & 6 bộ so sánh 5 bit; 4 NOT; 8 AND 2 ngõ vào; 2 splitter; 1 hằng số 0\\
\sig{HazardUnit} & 3 bộ so sánh 5 bit; 1 NOT; 1 OR; 1 AND 3 ngõ vào; 1 hằng số 0\\
\sig{main} & thanh ghi PC; ROM 256$\times$32; 2 bộ cộng; 5 MUX; 1 AND 32 bit; 1 OR; 4 hằng số; 8 sub-circuit; 5 probe\\
\sig{Pipeline} & 24 thanh ghi; ROM; 3 bộ cộng; 10 MUX; 2 NOT; 2 OR; 1 AND; 8 hằng số; 8 splitter; 10 sub-circuit; 8 probe\\
\bottomrule
\end{tabularx}
\end{center}
ROM lệnh của cả hai CPU đang nạp sẵn \texttt{demo.s} (13 từ, từ \texttt{00000513} đến \texttt{00008067}). Muốn chạy \texttt{hazards\_demo.s} phải nạp file \texttt{.hex} khác vào ROM.

\subsection{Phương trình đầy đủ của \texorpdfstring{\texttt{MainControl}}{MainControl}}
Lần theo dây vào từng cổng OR cho kết quả sau (bảy lớp lệnh ghi \sig{rd} được gom bằng hai cổng OR trung gian \sig{rw_a}, \sig{rw_b}):
\begin{align*}
\mathtt{RegWrite} &= \underbrace{(\mathtt{IsR}\lor\mathtt{IsI}\lor\mathtt{IsLoad})}_{\mathtt{rw\_a}}\lor\underbrace{(\mathtt{IsLui}\lor\mathtt{IsAuipc}\lor\mathtt{IsJAL})}_{\mathtt{rw\_b}}\lor\mathtt{IsJALR}\\
\mathtt{ALUSrc} &= \underbrace{(\mathtt{IsI}\lor\mathtt{IsLoad}\lor\mathtt{IsStore})}_{\mathtt{as\_a}}\lor\underbrace{(\mathtt{IsLui}\lor\mathtt{IsAuipc}\lor\mathtt{IsJALR})}_{\mathtt{as\_b}}\\
\mathtt{MemRead} &= \mathtt{IsLoad},\qquad \mathtt{MemWrite}=\mathtt{IsStore}\quad(\text{nối thẳng, không qua cổng})
\end{align*}
\begin{itemize}
  \item \textbf{LUI} có $\mathtt{ALUSrc}=1$ và $\mathtt{ASel}=10$: ALU tính $0+\mathtt{imm}$, nên kết quả chính là hằng số đã dịch trái 12 bit.
  \item \textbf{JAL} có $\mathtt{ALUSrc}=0$: ALU vẫn tính \sig{rs1+rs2} nhưng kết quả bị bỏ, vì \sig{WBSel}=10 ghi PC+4 và đích nhảy lấy từ bộ cộng \sig{PC+imm}.
  \item Bộ so sánh trong \sig{MainControl} đặt chế độ \emph{unsigned}, nhưng chỉ dùng ngõ ra ``bằng'', nên chế độ không ảnh hưởng.
\end{itemize}

\subsection{Bảng mã trong \texorpdfstring{\texttt{ALUControl}}{ALUControl}}
\begin{center}
\small
\begin{tabular}{@{}cll@{}}
\toprule
\sig{funct3} & \textbf{MUX kiểu R} & \textbf{MUX kiểu I}\\
\midrule
000 & MUX ADD/SUB (theo \sig{funct7[5]}) & hằng 0 (ADDI)\\
001 & 5 (SLL) & 5 (SLLI)\\
010 & 8 (SLT) & 8 (SLTI)\\
011 & 9 (SLTU) & 9 (SLTIU)\\
100 & 4 (XOR) & 4 (XORI)\\
101 & MUX SRL/SRA (dùng chung) & MUX SRL/SRA (dùng chung)\\
110 & 3 (OR) & 3 (ORI)\\
111 & 2 (AND) & 2 (ANDI)\\
\bottomrule
\end{tabular}
\end{center}
\begin{itemize}
  \item MUX SRL/SRA có một ngõ ra nối vào \emph{cả hai} bảng, vì \texttt{srai} và \texttt{sra} đều phân biệt bằng bit 30 của lệnh (\sig{funct7[5]}).
  \item Kiểu I với \sig{funct3}=000 dùng thẳng hằng số 0, \emph{không} qua MUX ADD/SUB. Ở \texttt{addi} bit 30 là bit 10 của hằng số, nên \texttt{addi x1,x0,-1024} (bit 30 = 1) vẫn là phép cộng. Đây là chỗ dễ sai nhất của bảng này.
  \item MUX cuối chọn theo \sig{ALUOp}: ngõ 0 hằng 0 (ép ADD), ngõ 1 hằng 1 (ép SUB), ngõ 2 bảng R, ngõ 3 bảng I.
\end{itemize}

\subsection{Mười multiplexer của \texorpdfstring{\texttt{Pipeline}}{Pipeline}}
\begin{center}
\small
\begin{tabularx}{\linewidth}{@{}l>{\raggedright}p{2.6cm}>{\raggedright\arraybackslash}X@{}}
\toprule
\textbf{Ra} & \textbf{Chọn theo} & \textbf{Các ngõ vào (0, 1, 2, 3)}\\
\midrule
\sig{next_pc} & \sig{pc_src_ex} & \sig{pc_plus4_if}, \sig{target_ex}\\
\sig{instr_in} & \sig{pc_src_ex} & lệnh từ ROM, 0 (flush lệnh vừa nạp)\\
\sig{ctrl_in} (13 bit) & \sig{bubble} & \sig{ctrl_id}, 0 (bubble)\\
\sig{fwd_mem_val} & \sig{WBSel1_mem} & \sig{alu_result_mem}, \sig{pc_plus4_mem}\\
\sig{opA} & \sig{ForwardA} & \sig{rs1_data_ex}, \sig{wb_data}, \sig{fwd_mem_val}, \sig{rs1_data_ex}\\
\sig{opB} & \sig{ForwardB} & \sig{rs2_data_ex}, \sig{wb_data}, \sig{fwd_mem_val}, \sig{rs2_data_ex}\\
\sig{alu_a} & \sig{ASel_ex} & \sig{opA}, \sig{pc_ex}, 0, 0\\
\sig{alu_b} & \sig{ALUSrc_ex} & \sig{opB}, \sig{imm_ex}\\
\sig{target_ex} & \sig{IsJALR_ex} & \sig{pc_target} (\sig{pc_ex+imm}), \sig{jalr_target} (ALU AND \texttt{0xFFFFFFFE})\\
\sig{wb_data} & \sig{WBSel_wb} & \sig{alu_result_wb}, \sig{mem_rdata_wb}, \sig{pc_plus4_wb}, 0\\
\bottomrule
\end{tabularx}
\end{center}
\begin{itemize}
  \item \textbf{Mã không dùng vẫn an toàn:} ngõ 3 của MUX forwarding quay về giá trị register file; ngõ 3 của \sig{ASel} và \sig{WBSel} là 0.
  \item \textbf{Forwarding nằm trước ALUSrc:} \sig{opB} được forward trước rồi mới chọn giữa \sig{opB} và \sig{imm_ex}. Lệnh kiểu I có thể bật \sig{ForwardB} vì trường \sig{rs2} thực chất là bit hằng số, nhưng vô hại vì \sig{ALUSrc}=1 bỏ qua \sig{opB}. Còn \sig{opB} thì đồng thời là \sig{store_data_mem}, nên lệnh store nhận được dữ liệu đã forward.
  \item \textbf{BranchUnit và JALR dùng toán hạng đã forward:} \sig{BranchUnit} nhận \sig{opA}/\sig{opB}, còn đích JALR lấy từ ALU. Vì vậy \texttt{beq} hay \texttt{ret} đứng ngay sau lệnh ALU ghi thanh ghi nguồn đều không phải stall.
  \item \textbf{PC+4 được tính lại ở EX:} ID/EX không chứa PC+4. Bộ cộng thứ ba tính \sig{pc_ex+4} rồi đưa vào \sig{pc_plus4_mem} làm địa chỉ trả về của JAL/JALR.
  \item \textbf{\sig{bubble} = \sig{stall} OR \sig{pc_src_ex}.} Hai tín hiệu không bao giờ cùng bằng 1, vì \sig{stall} cần một lệnh load ở EX còn \sig{pc_src_ex} cần một lệnh nhảy ở EX, mà EX chỉ chứa một lệnh.
\end{itemize}

\subsection{Đóng gói từ điều khiển}
Splitter 13 bit ở ID và EX chia từ điều khiển thành 10 trường theo thứ tự bit 0 $\rightarrow$ 12:
\sig{RegWrite}(0), \sig{ALUSrc}(1), \sig{MemWrite}(2), \sig{MemRead}(3), \sig{IsBranch}(4), \sig{ALUOp}(5--6), \sig{ASel}(7--8), \sig{WBSel}(9--10), \sig{IsJAL}(11), \sig{IsJALR}(12).
Sang EX/MEM còn 5 bit (\sig{RegWrite}, \sig{MemWrite}, \sig{MemRead}, \sig{WBSel}); sang MEM/WB còn 3 bit (\sig{RegWrite}, \sig{WBSel}).

\subsection{Xung nhịp}
\begin{itemize}
  \item PC, 23 thanh ghi pipeline và RAM dữ liệu ghi ở \textbf{sườn lên} của \sig{CLK}.
  \item \sig{RegisterFile} trong \sig{Pipeline} nhận \sig{CLK_n} (qua cổng NOT) nên ghi ở \textbf{sườn xuống}, giữa chu kỳ.
  \item Trong \sig{main} chỉ có một miền xung nhịp: register file, PC và RAM cùng ghi ở sườn lên.
\end{itemize}

\subsection{Những điểm mạch \emph{không} làm (nên biết khi bị hỏi)}
\begin{itemize}
  \item \textbf{\sig{MemRead} không điều khiển bộ nhớ.} \sig{Memory} không có chân \sig{MemRead}; RAM luôn đọc bất đồng bộ. Trong \sig{main}, tunnel \sig{MemRead} chỉ có một đầu (không nối đi đâu). Trong \sig{Pipeline}, chỉ \sig{MemRead_ex} được dùng (vào \sig{HazardUnit}); \sig{MemRead_mem} được mang theo nhưng không có nơi nhận.
  \item \textbf{\sig{HazardUnit} stall thừa.} Mạch so \sig{rd_ex} với \sig{rs1_id} và \sig{rs2_id} mà không xét lệnh có thật sự dùng hai thanh ghi đó không. Ví dụ \texttt{lw x2,0(x0)} rồi \texttt{addi x3,x5,2}: bit 24:20 của \texttt{addi} là hằng số 2, trùng \sig{rd_ex}=2, nên vẫn stall 1 chu kỳ dù không cần. Kết quả vẫn đúng, chỉ chậm hơn. Muốn sửa: thêm tín hiệu ``dùng \sig{rs2}'' (chỉ kiểu R, S, B) và ``dùng \sig{rs1}'' (mọi lệnh trừ LUI, AUIPC, JAL).
  \item \textbf{Không bẫy truy cập lệch (misaligned).} \texttt{lw} bỏ qua \sig{Address[1:0]} và đọc nguyên từ; \texttt{lh} chỉ xét \sig{Address[1]}. Trong ROM của \sig{StoreMerge} (\texttt{1 2 4 8 3 3 c c f f f f f f f f}), \texttt{sh} ở độ lệch 1 ghi như độ lệch 0, ở độ lệch 3 ghi như độ lệch 2; \sig{funct3} có size 11 được coi như ghi từ.
  \item \textbf{Địa chỉ bị lặp mỗi 1 KB.} RAM chỉ dùng \sig{Address[9:2]}, bit 31:10 bị bỏ, nên \texttt{0x400} trỏ về cùng từ 0 như \texttt{0x000}. ROM lệnh cũng chỉ dùng \sig{PC[9:2]}: PC vượt \texttt{0x3FC} sẽ quay về đầu chương trình.
  \item \textbf{Opcode lạ thành NOP.} Mọi cờ lớp lệnh bằng 0 nên không ghi thanh ghi, không ghi bộ nhớ, không nhảy. Mạch không báo lỗi lệnh không hợp lệ.
\end{itemize}

\begin{kiemtra}
\begin{itemize}
  \item Vì sao \texttt{addi x1,x0,-1024} không bị biến thành phép trừ?
  \item Trong \sig{Pipeline}, \texttt{addi x1,x0,5} rồi ngay sau là \texttt{beq x1,x0,L}: giá trị x1 đến \sig{BranchUnit} bằng đường nào?
  \item Viết một cặp lệnh làm \sig{HazardUnit} stall dù không có phụ thuộc thật.
  \item Nếu xoá tunnel \sig{MemRead} trong \sig{main}, mạch có chạy khác đi không?
\end{itemize}
\end{kiemtra}

% =====================================================================
\section{Demo từng bước}
\label{sec:demo}
% =====================================================================
\subsection{Chương trình \texttt{demo.s}}
\begin{center}
\small
\begin{tabular}{@{}llll@{}}
\toprule
\textbf{Địa chỉ} & \textbf{Mã máy} & \textbf{Lệnh} & \textbf{Ý nghĩa}\\
\midrule
0x00 & \texttt{00000513} & \texttt{li x10,0} & tổng = 0\\
0x04 & \texttt{00a00593} & \texttt{li x11,10} & biến đếm = 10\\
0x08 & \texttt{00b50533} & \texttt{loop: add x10,x10,x11} & tổng += biến đếm\\
0x0c & \texttt{fff58593} & \texttt{addi x11,x11,-1} & biến đếm $-=1$\\
0x10 & \texttt{fe059ce3} & \texttt{bnez x11,loop} & lặp nếu khác 0\\
0x14 & \texttt{00a02023} & \texttt{sw x10,0(x0)} & lưu tổng\\
0x18 & \texttt{00002603} & \texttt{lw x12,0(x0)} & nạp lại vào x12\\
0x1c & \texttt{010000ef} & \texttt{jal x1,double} & gọi hàm, x1 = 0x20\\
0x20 & \texttt{40a606b3} & \texttt{sub x13,x12,x10} & $110-55=55$\\
0x24 & \texttt{0646a713} & \texttt{slti x14,x13,100} & $55<100\Rightarrow1$\\
0x28 & \texttt{0000006f} & \texttt{halt: j halt} & vòng lặp kết thúc\\
0x2c & \texttt{00c60633} & \texttt{double: add x12,x12,x12} & x12 = 110\\
0x30 & \texttt{00008067} & \texttt{ret} & về địa chỉ trong x1\\
\bottomrule
\end{tabular}
\end{center}
Kết quả: tổng $10+9+\dots+1=55=\mathtt{0x37}$; sau lời gọi $x12=110=\mathtt{0x6e}$; $x13=\mathtt{0x37}$; $x14=1$; từ dữ liệu 0 $=\mathtt{0x37}$; PC cuối $=\mathtt{0x28}$; $x1=\mathtt{0x20}$.

\subsection{Cách đọc một tick}
\begin{itemize}
  \item Một lần Ctrl+F9 = một chu kỳ xung nhịp = một lệnh (single-cycle).
  \item Sau tick, PC đã trỏ tới lệnh \emph{kế tiếp}. Hàng probe hiện lệnh kế tiếp (\sig{instr}) và nó \emph{sẽ} tính ra gì; tab State hiện các thanh ghi mà lệnh \emph{trước} vừa ghi.
  \item Nhịp độ: nói một câu, bấm một tick, chỉ một thứ, chờ hai giây.
\end{itemize}

\subsection{Một lệnh đi qua CPU, theo từng nhóm}
\begin{description}[style=nextline,leftmargin=1.2em,labelindent=0pt]
  \item[Kiểu R: \texttt{add x10,x10,x11} (tick 3)] ID: $\mathtt{ALUOp}=10$, $\mathtt{ALUSrc}=0$; EX: $\mathtt{ALUCtrl}=0$, kết quả \texttt{0x0a}; WB: ghi vào $x10$.
  \item[Kiểu I: \texttt{addi x11,x11,-1} (tick 4)] hằng số \texttt{0xffffffff}; $\mathtt{ALUSrc}=1$; $\mathtt{ALUOp}=11$; kết quả \texttt{0x09}.
  \item[Branch rẽ: \texttt{bnez x11,loop} (tick 5)] hằng số kiểu B $=-8$; \sig{BranchUnit} so $x11=9$ với $x0=0$ bằng $\mathtt{funct3}=001$ (BNE): $\mathtt{branch\_taken}=1$; $\mathtt{pc\_target}=\mathtt{0x10}-8=\mathtt{0x08}$.
  \item[Store: \texttt{sw x10,0(x0)} (tick 33)] $\mathtt{MemWrite}=1$; địa chỉ $0$; \sig{StoreMerge} bật cả bốn byte; từ 0 của RAM thành \texttt{0x37}.
  \item[Load: \texttt{lw x12,0(x0)} (tick 34)] \sig{LoadExtract} cho qua nguyên từ; $\mathtt{WBSel}=01$; $x12=\mathtt{0x37}$.
  \item[JAL: \texttt{jal x1,double} (tick 35)] $\mathtt{pc\_target}=\mathtt{0x1c}+\mathtt{0x10}=\mathtt{0x2c}$; $\mathtt{WBSel}=10$: $x1=\mathtt{0x20}$.
  \item[JALR: \texttt{ret} (tick 37)] ALU tính $x1+0=\mathtt{0x20}$; xóa bit 0; $\mathtt{next\_pc}=\mathtt{0x20}$; $\mathtt{rd}=x0$ nên không ghi gì.
  \item[SUB và SLTI (tick 38, 39)] $\mathtt{funct7}=0100000$ chọn SUB: $\mathtt{0x6e}-\mathtt{0x37}=\mathtt{0x37}$; \texttt{slti}: $\mathtt{0x37}<100\Rightarrow x14=1$.
  \item[Kết thúc: \texttt{halt: j halt}] nhảy về chính nó, PC đứng ở \texttt{0x28}.
\end{description}

\subsection{Cả 40 tick của \texttt{demo.s} (single-cycle)}
Tick $n$ thực thi lệnh ở PC \emph{trước} tick; hai cột cuối là trạng thái \emph{sau} tick.

\small
\begin{longtable}{@{}rllll@{}}
\toprule
\textbf{Tick} & \textbf{PC chạy} & \textbf{Lệnh} & \textbf{Thay đổi} & \textbf{PC sau}\\
\midrule
\endhead
1 & 0x00 & \texttt{li x10,0} & (x10 vẫn 0) & 0x04\\
2 & 0x04 & \texttt{li x11,10} & x11 = 0x0a & 0x08\\
3 & 0x08 & \texttt{add} & x10 = 0x0a & 0x0c\\
4 & 0x0c & \texttt{addi} & x11 = 0x09 & 0x10\\
5 & 0x10 & \texttt{bnez} (rẽ) & -- & 0x08\\
6 & 0x08 & \texttt{add} & x10 = 0x13 & 0x0c\\
7 & 0x0c & \texttt{addi} & x11 = 0x08 & 0x10\\
8 & 0x10 & \texttt{bnez} (rẽ) & -- & 0x08\\
9 & 0x08 & \texttt{add} & x10 = 0x1b & 0x0c\\
10 & 0x0c & \texttt{addi} & x11 = 0x07 & 0x10\\
11 & 0x10 & \texttt{bnez} (rẽ) & -- & 0x08\\
12 & 0x08 & \texttt{add} & x10 = 0x22 & 0x0c\\
13 & 0x0c & \texttt{addi} & x11 = 0x06 & 0x10\\
14 & 0x10 & \texttt{bnez} (rẽ) & -- & 0x08\\
15 & 0x08 & \texttt{add} & x10 = 0x28 & 0x0c\\
16 & 0x0c & \texttt{addi} & x11 = 0x05 & 0x10\\
17 & 0x10 & \texttt{bnez} (rẽ) & -- & 0x08\\
18 & 0x08 & \texttt{add} & x10 = 0x2d & 0x0c\\
19 & 0x0c & \texttt{addi} & x11 = 0x04 & 0x10\\
20 & 0x10 & \texttt{bnez} (rẽ) & -- & 0x08\\
21 & 0x08 & \texttt{add} & x10 = 0x31 & 0x0c\\
22 & 0x0c & \texttt{addi} & x11 = 0x03 & 0x10\\
23 & 0x10 & \texttt{bnez} (rẽ) & -- & 0x08\\
24 & 0x08 & \texttt{add} & x10 = 0x34 & 0x0c\\
25 & 0x0c & \texttt{addi} & x11 = 0x02 & 0x10\\
26 & 0x10 & \texttt{bnez} (rẽ) & -- & 0x08\\
27 & 0x08 & \texttt{add} & x10 = 0x36 & 0x0c\\
28 & 0x0c & \texttt{addi} & x11 = 0x01 & 0x10\\
29 & 0x10 & \texttt{bnez} (rẽ) & -- & 0x08\\
30 & 0x08 & \texttt{add} & x10 = 0x37 (= 55) & 0x0c\\
31 & 0x0c & \texttt{addi} & x11 = 0x00 & 0x10\\
32 & 0x10 & \texttt{bnez} (\textbf{không} rẽ) & -- & 0x14\\
33 & 0x14 & \texttt{sw} & từ 0 = 0x37 & 0x18\\
34 & 0x18 & \texttt{lw} & x12 = 0x37 & 0x1c\\
35 & 0x1c & \texttt{jal} & x1 = 0x20 & \textbf{0x2c}\\
36 & 0x2c & \texttt{add x12,x12,x12} & x12 = 0x6e (= 110) & 0x30\\
37 & 0x30 & \texttt{ret} & -- & \textbf{0x20}\\
38 & 0x20 & \texttt{sub} & x13 = 0x37 & 0x24\\
39 & 0x24 & \texttt{slti} & x14 = 0x01 & 0x28\\
40 & 0x28 & \texttt{halt} & -- & 0x28 (đứng yên)\\
\bottomrule
\end{longtable}
\normalsize

\textbf{Vì sao kết quả ghi 39 chu kỳ chứ không phải 40:} tick 40 chỉ lặp lại lệnh nhảy cuối; lần ghi thanh ghi cuối cùng xảy ra ở tick 39. \textbf{Các điểm nên chỉ:} PC \emph{quay lùi} ở tick 5, 8, 11, \dots{} (vòng lặp); ở tick 32 nó \emph{đi tiếp}; ở tick 35 nó \emph{nhảy tới} \texttt{0x2c}; ở tick 37 nó \emph{trở về} \texttt{0x20}.

\subsection{Các giá trị probe cần xem (single-cycle)}
Giá trị hiện ở hàng probe \emph{sau} tick, tức của lệnh ở PC mới:
\begin{center}
\small
\begin{tabular}{@{}rlllll@{}}
\toprule
\textbf{Sau tick} & \textbf{PC} & \sig{instr} & \sig{alu_result} & \sig{branch_taken} & \sig{wb_data}\\
\midrule
1 & 04 & \texttt{00a00593} & 0a & 0 & 0a\\
2 & 08 & \texttt{00b50533} & 0a & 0 & 0a\\
3 & 0c & \texttt{fff58593} & 09 & 0 & 09\\
4 & 10 & \texttt{fe059ce3} & 09 & \textbf{1} & 09\\
5 & \textbf{08} & \texttt{00b50533} & 13 & 0 & 13\\
31 & 10 & \texttt{fe059ce3} & 00 & \textbf{0} & 00\\
\bottomrule
\end{tabular}
\end{center}

% =====================================================================
\section{Chia việc cho nhóm ba người}
\label{sec:chia}
% =====================================================================
\begin{center}
\small
\begin{tabularx}{\linewidth}{@{}c>{\raggedright}p{4.2cm}>{\raggedright}X>{\raggedright\arraybackslash}p{2.4cm}@{}}
\toprule
\textbf{Người} & \textbf{Chủ đề} & \textbf{Circuit} & \textbf{Bước demo}\\
\midrule
A & front end, giải mã, single-cycle & \sig{main}, \sig{InstructionDecode}, \sig{MainControl}, \sig{ImmGen}, \sig{RegisterFile} & 1--5\\
B & thực thi và bộ nhớ, kiểm thử & \sig{ALU}, \sig{ALUControl}, \sig{BranchUnit}, \sig{Memory}, \sig{LoadExtract}, \sig{StoreMerge} & 6, 9, 11, 12\\
C & pipeline và hazard & \sig{Pipeline}, \sig{ForwardingUnit}, \sig{HazardUnit} & 7, 8, 10\\
\bottomrule
\end{tabularx}
\end{center}
Cả 14 circuit đều có người phụ trách: A có 5, B có 6, C có 3 (nhưng \sig{Pipeline} là circuit lớn nhất).

% =====================================================================
\section{Tự đánh giá}
% =====================================================================
\begin{center}
\begin{tabular}{@{}cp{11cm}@{}}
\toprule
\textbf{Mức} & \textbf{Bạn có thể \dots}\\
\midrule
1 & gọi tên mọi circuit và nói nó làm gì trong một câu\\
2 & liệt kê các linh kiện bên trong mọi circuit và nói vì sao có chúng\\
3 & lần theo bất kỳ lệnh nào qua \sig{main} với giá trị thật\\
4 & dự đoán hành vi pipeline (stall, forwarding, flush) của một chương trình ngắn\\
5 & giải thích dự án được kiểm thử thế nào và vì sao tin được kết quả\\
\bottomrule
\end{tabular}
\end{center}
Mục tiêu: mức 5 cho phần của riêng bạn và mức 3 cho hai phần còn lại.

\textbf{Bài thi cuối (không xem tài liệu):} (1) vẽ \sig{main} và \sig{Pipeline}; (2) lần theo \texttt{lw} theo sau bởi một \texttt{addi} phụ thuộc qua pipeline; (3) giải thích ba lỗi đã gặp; (4) giải thích vì sao có thể tin phần chứng minh.

\textbf{Thứ tự đọc gợi ý:} \texttt{README.md} $\rightarrow$ \texttt{docs/GUIDE.md} $\rightarrow$ \texttt{docs/REPORT.md} $\rightarrow$ \texttt{docs/COMPONENTS.md} $\rightarrow$ \texttt{docs/DEMO\_EXPLAIN.md} $\rightarrow$ \texttt{docs/PRESENTATION.md} $\rightarrow$ \texttt{docs/TOOLS.md}, \texttt{docs/DESIGN\_LOG.md}.

\end{document}
